\chapter{One complete session}
\label{cha:worked-example}

This appendix walks through one real run of the pipeline described in \Cref{cha:design}: the plan the model committed to and the totals, caveats and rationale computed from it, alongside a food-database lookup that illustrates the tool the plan is built on. The run is the \emph{regular} scenario (\Cref{sec:methodology-scenarios}: a fixed set of macro targets -- roughly 2728~kcal, 124.8~g protein, 386.7~g carbohydrate, 75.8~g fat -- and no free-form request) planned by \texttt{gemini-3.7-flash} with the totaller and reflection loop both active. It is one illustrative run, not a measurement; the systematic figures belong to \Cref{cha:results}.

\Cref{lst:example-lookup} shows a representative response of the lookup tool -- a single query for \texttt{rolled oats} against the same catalogue, capped at one hit per source. The response is structured JSON returning both catalogue channels at once: a branded, Polish-language Open Food Facts entry and a curated, generic USDA one, each carrying the per-100~g nutrients and the source-prefixed \texttt{code} the agent must copy verbatim.

\lstinputlisting[caption={The food-database lookup tool's JSON response for one query, one hit per source. The agent copies \texttt{code} and \texttt{name} from a hit and may read its nutrient values, but never emits nutrient numbers of its own.},label={lst:example-lookup}]{listings/example_lookup.json}

\Cref{lst:example-plan} shows the breakfast meal exactly as the model emitted it. Each portion is a reference -- a \texttt{code}, the \texttt{name} copied from the lookup, and a mass in grams -- and nothing more. The single meal freely mixes both catalogues -- Open Food Facts codes for the Polish branded oats, milk and honey, USDA codes for the generic chia, apple, walnuts and cinnamon.

\lstinputlisting[caption={The breakfast meal as the nutrition agent emitted it: references only (code, name, grams) plus preparation instructions, no nutrient values.},label={lst:example-plan}]{listings/example_plan.json}

Only after the plan is hydrated -- every reference resolved back to a real database row and the totaller run over the result -- does the deterministic layer fill in the numbers. \Cref{tab:example-breakfast-items} shows the same breakfast meal at that next step. The output carries an equivalent breakdown for every meal in the plan; only one is shown here.

\begin{table}[htbp]
    \centering
    \footnotesize
    \setlength{\tabcolsep}{3pt}
    \begin{tabular}{
    l
    S[table-format=3.0]
    S[table-format=3.0]
    S[table-format=2.1]
    S[table-format=3.1]
    S[table-format=2.1]
}
\toprule
{Product} & {Mass [\si{\gram}]} & {Energy [\si{\kcal}]} & {Protein [\si{\gram}]} & {Carbohydrate [\si{\gram}]} & {Fat [\si{\gram}]} \\
\midrule
Płatki owsiane górskie & 90 & 337 & 11.7 & 54.9 & 6.0 \\
Seeds, chia seeds, dried & 20 & 97 & 3.3 & 8.4 & 6.1 \\
Mleko 2\% & 250 & 125 & 8.0 & 11.7 & 5.0 \\
Apples, gala, with skin, raw & 150 & 92 & 0.2 & 22.2 & 0.2 \\
Nuts, walnuts, English, halves, raw & 27 & 197 & 3.9 & 2.9 & 18.8 \\
Miód nektarowy wielokwiatowy & 10 & 34 & 0.1 & 8.3 & 0.0 \\
Spices, cinnamon, ground & 2 & 5 & 0.1 & 1.6 & 0.0 \\
\midrule
\textbf{Meal total} & {--} & 886 & 27.3 & 110.1 & 36.3 \\
\multicolumn{6}{p{0.92\linewidth}}{\textit{\qty{20.6}{\gram} fibre, \qty{0.26}{\gram} salt, \qty{107}{\milli\gram} sodium, \qty{702.67}{\milli\gram} potassium, \qty{518.94}{\milli\gram} calcium, \qty{6.05}{\milli\gram} iron, \qty{0.4}{\milli\gram} vitamin~C}} \\
\bottomrule
\end{tabular}

    \caption{The breakfast meal after hydration: per-portion masses and macronutrients, meal total, and additional micronutrients. Long product names are trimmed; the output carries the same structure for all four meals.}
    \label{tab:example-breakfast-items}
\end{table}

\Cref{tab:example-mealplan} is the central artefact of the output: the whole plan laid out one row per meal, each with its products, meal totals and the preparation recipe. Product lists and recipes are trimmed here to fit the page.

\begin{table}[htbp]
    \centering
    \scriptsize
    \setlength{\tabcolsep}{4pt}
    \renewcommand{\arraystretch}{1.15}
    \begin{tabularx}{\textwidth}{@{}l >{\raggedright\arraybackslash}p{2.1cm} >{\raggedright\arraybackslash}p{3.4cm} >{\raggedright\arraybackslash}p{1.9cm} X@{}}
\toprule
Meal & Name & Products & Meal totals & Recipe \\
\midrule
breakfast & Overnight oats with chia [\ldots] &
Płatki owsiane górskie: \qty{90}{\gram} \newline Seeds, chia seeds, dried: \qty{20}{\gram} \newline Mleko 2\%: \qty{250}{\gram} \newline Apples, gala, with skin, raw: \qty{150}{\gram} \newline (+3 more) &
\qty{886}{\kcal} \newline P27.3 C110.1 F36.3 &
In a bowl or mason jar, combine the rolled oats, chia seeds, and ground cinnamon. [\ldots] \\
\addlinespace
lunch & Chicken and vegetable stir-fry [\ldots] &
Rice stick noodle: \qty{130}{\gram} \newline Chicken, breast, boneless, skinless, raw: \qty{130}{\gram} \newline Broccoli, raw: \qty{150}{\gram} \newline Peppers, sweet, red, raw: \qty{100}{\gram} \newline (+5 more) &
\qty{854}{\kcal} \newline P46.7 C135.0 F16.5 &
Prepare the ingredients: cut chicken breast into bite-sized strips; cut broccoli into small florets; slice red bell pepper and carrot into thin matchsticks; finely mince garlic and fresh ginger root. [\ldots] \\
\addlinespace
dinner & Pan-seared cod fillet with [\ldots] &
Fish, cod, Atlantic, raw: \qty{130}{\gram} \newline Potatoes, flesh and skin, raw: \qty{460}{\gram} \newline Beans, snap, green, raw: \qty{200}{\gram} \newline Dill weed, fresh: \qty{10}{\gram} \newline (+2 more) &
\qty{714}{\kcal} \newline P36.7 C95.1 F22.9 &
Thoroughly scrub and rinse the potatoes, cut them into bite-sized pieces, and place them in a saucepan of boiling water for 15-20 minutes until tender; then drain. [\ldots] \\
\addlinespace
snack & Natural skyr with fresh [\ldots] &
Skyr naturalny: \qty{100}{\gram} \newline Blueberries, raw: \qty{120}{\gram} \newline Bananas, raw: \qty{120}{\gram} \newline Miód nektarowy wielokwiatowy: \qty{10}{\gram} &
\qty{273}{\kcal} \newline P14.2 C57.2 F0.8 &
Spoon the natural skyr into a serving bowl. [\ldots] \\
\bottomrule
\end{tabularx}

    \caption{The full meal plan in the run's output: one row per meal, with products (mass per portion), the meal's energy and macronutrient totals, and the recipe. Product lists are trimmed to the first few items and recipes to the first step (\texttt{[\ldots]}).}
    \label{tab:example-mealplan}
\end{table}

At whole-plan level the output also carries the day's nutrient totals against the run's macro targets (\Cref{tab:example-totals}), the data-coverage caveats, and a patient-facing rationale. Energy, protein and fat land within about a percent of target and carbohydrate within a few percent (the reflection loop reported no issues, so it ran zero refinement passes here). Fiber sits well above its target; it is reported for information but, being the nutrient most often missing from a crowd-sourced record, is kept out of the macro error metrics rather than treated as an accuracy failure (\Cref{sec:impl-missing-nutrients}).

\begin{table}[htbp]
    \centering
    \footnotesize
    \begin{tabular}{lrr}
\toprule
Nutrient & Actual & Target \\
\midrule
Energy [kcal]      & 2727  & 2728 \\
Protein [g]        & 124.9 & 124.8 \\
Carbohydrate [g]   & 397.5 & 386.7 \\
Fat [g]            & 76.5  & 75.8 \\
Saturated fat [g]  & 17.0  & -- \\
Fiber [g]          & 55.0  & 25.0 \\
Sugar [g]          & 98.2  & -- \\
Salt [g]           & 0.4   & -- \\
Sodium [mg]        & 1205  & -- \\
Potassium [mg]     & 6095  & -- \\
Calcium [mg]       & 851   & -- \\
Iron [mg]          & 17    & -- \\
Vitamin C [mg]     & 417   & -- \\
Vitamin D [$\mu$g] & 1.3   & -- \\
Cholesterol [mg]   & 172   & -- \\
\bottomrule
\end{tabular}

    \caption{The day's nutrient totals against the run's macro targets, computed after hydration and totalling. A dash marks a nutrient the run does not target.}
    \label{tab:example-totals}
\end{table}

The same totalling step also reports which totals are understated because a food leaves a nutrient blank (\Cref{tab:example-caveats}). Salt is flagged across the twenty USDA foods that do not model it, while the handful of Open Food Facts entries lacking a micronutrient value (the milk, honey and skyr) are flagged for the nutrients they genuinely leave blank.

\begin{table}[htbp]
    \centering
    \footnotesize
    \begin{tabular}{lrrr}
\toprule
Nutrient & Foods without data & Plan mass affected & Reported total \\
\midrule
Salt          & 20 of 26 & 75\% & 0.38 g \\
Cholesterol   & 4 of 26  & 16\% & 171.95 mg \\
Fiber         & 4 of 26  & 16\% & 55.04 g \\
Iron          & 4 of 26  & 16\% & 17.36 mg \\
Potassium     & 4 of 26  & 16\% & 6094.88 mg \\
Vitamin C     & 4 of 26  & 16\% & 417.17 mg \\
Vitamin D     & 4 of 26  & 16\% & 1.33 $\mu$g \\
Calcium       & 3 of 26  & 5\%  & 851.05 mg \\
Saturated fat & 2 of 26  & 1\%  & 16.96 g \\
\bottomrule
\end{tabular}

    \caption{Data-coverage caveats reported alongside the totals: for each nutrient, how many plan foods carry no value for it, the share of plan mass they account for, and the reported total that is therefore understated.}
    \label{tab:example-caveats}
\end{table}

The output also carries the plan's \emph{rationale} field. The excerpt in \Cref{fig:example-rationale} shows the shape of that text: an opening summary, a list of dietary choices, and a supplementation section. The latter is the kind of note the nutrition agent's prompt was extended to invite (\Cref{sec:impl-prompt-fixes}): here it flags vitamin~D as under-covered by food alone, citing the day's reported total of about 1.3~$\mu$g.

\begin{figure}[htbp]
    \centering
    \footnotesize
    \begin{minipage}{0.92\linewidth}
        \textbf{Rationale:}\\[0.4em]
        This 1-day meal plan is designed for a 30-year-old male omnivore aiming for \textasciitilde{}2,728 kcal with balanced macronutrients (124.8 g protein, 386.7 g carbohydrates, and 75.8 g fat).

Key dietary choices:
- High-quality, lean protein sources such as chicken breast, Atlantic cod, natural skyr, and semi-skimmed milk provide complete amino acid profiles to comfortably hit the 124.8 g daily protein target while keeping saturated fats well below 10\% of total energy.
[...]

Supplementation guidance:
- Vitamin D: Like most adults residing in temperate European latitudes with limited year-round sun exposure, dietary sources alone (providing \textasciitilde{}1.3 µg) are insufficient to reach optimal physiological levels (15–20 µg / 600–800 IU daily). Supplementation with Vitamin D3 (cholecalciferol) at 1,000–2,000 IU (25–50 µg) daily is recommended, especially during autumn and winter months; confirm specific dosing with a healthcare professional based on 25(OH)D blood testing.
[...]


    \end{minipage}
    \caption{Opening of the patient-facing rationale in the run's output. Only the start of each section is kept; the rest is trimmed with \texttt{[\ldots]}.}
    \label{fig:example-rationale}
\end{figure}
